\begin{abstract}
Existing Korean benchmarks for Large Language Models (LLMs) evaluate performance across diverse knowledge domains, but their accuracy scores cannot distinguish whether incorrect responses arise from a lack of required knowledge or a failure to apply that knowledge correctly.
Our analysis of existing Korean benchmarks reveals that most questions can be answered through factual recall alone, while those demanding greater application also differ in required knowledge, making it difficult to isolate application ability.
To address this, we introduce \name{}, a framework that constructs \textit{knowledge-matched contrasts} by varying application demand while preserving the required knowledge.
During construction, \name{} identifies the knowledge underlying each source question as an explicit anchor bundle to guide generation of questions requiring more complex application.
Generated candidates then undergo automated and human validation, with iterative revision based on evaluation feedback. 
Experiments demonstrate that the resulting knowledge-matched sets largely satisfy the intended constraints, enabling controlled evaluation across diverse LLMs.
These evaluations reveal that increasing application demand exposes failures that persist even when the required knowledge is provided.
\end{abstract}
